% Generated draft: requires mathematical and source review.
\documentclass[11pt]{article}
\usepackage[margin=25mm]{geometry}
\usepackage{fontspec}
\setmainfont{Latin Modern Roman}
\usepackage{amsmath,amssymb,mathtools,unicode-math}
\setmathfont{Latin Modern Math}
\usepackage{enumitem,needspace,xurl,hyperref}
\setlength{\parindent}{0pt}
\setlength{\parskip}{5pt}
\setlength{\emergencystretch}{3em}
\newcommand{\N}{\mathbb{N}}
\newcommand{\Z}{\mathbb{Z}}
\newcommand{\Q}{\mathbb{Q}}
\newcommand{\R}{\mathbb{R}}
\newcommand{\C}{\mathbb{C}}
\newcommand{\F}{\mathbb{F}}
\newcommand{\A}{\mathbb{A}}
\newcommand{\PP}{\mathbb{P}}
\newcommand{\E}{\mathbb{E}}
\newcommand{\eps}{\varepsilon}
\newcommand{\dd}{\,\mathrm{d}}
\begin{document}
\section*{UnsolvedMath ID: 30003295}
\textbf{Source problem code:} catalogue entry
\subsection{Definitions and mathematical statement}
Let \(q=p^r\), where \(p\) is prime and \(r\geq 1\) is an integer, and let \(\F_q\) be the field with \(q\) elements. For each integer \(g\geq 2\), interpret \(\mathcal M_g\), explicitly called a coarse moduli space in the supplied record, as the coarse moduli space over \(\F_q\) of smooth, proper, geometrically connected, unmarked curves of genus \(g\). This uses the standard interpretation of the notation; the supplied statement does not explicitly specify these curve hypotheses. No compactification or boundary of singular stable curves is included under this interpretation.

Define
\[
N_g(q)=\#\mathcal M_g(\F_q),
\qquad
\mathcal M_g(\F_q)=\operatorname{Hom}_{\F_q}\bigl(\operatorname{Spec}\F_q,\mathcal M_g\bigr).
\]
Thus \(N_g(q)\) counts rational points of the coarse moduli space. It is not defined here as a count of \(\F_q\)-isomorphism classes of curves or as a sum weighted by reciprocals of automorphism-group orders.

For each fixed prime power \(q\), does there exist a real number \(L_q\) with \(0<L_q<\infty\) such that
\[
\lim_{\substack{g\to\infty\\g\in\Z,\ g\geq 2}}
\frac{N_g(q)}{q^{3g-3}}=L_q\,?
\]
The question concerns the limit through all integer genera, with \(q\) held fixed. The curve-moduli interpretation above requires confirmation against the original problem discussion.
\subsection{Short English statement}
For a fixed finite field \(\F_q\), count its rational points on the coarse moduli space of smooth genus-\(g\) curves. Does dividing this count by \(q^{3g-3}\) give a sequence converging to a finite, strictly positive number as \(g\) tends to infinity? The standard curve-moduli interpretation is implicit in the supplied source and requires review.
\subsection{Sources}
\begin{itemize}
\item Supplied local UnsolvedMath registry record, permanent ID \(30003295\), title ``Asymptotic Point Counts on Curve Moduli Spaces,'' exact source problem code \texttt{catalogue entry}, in the collection ``Oberwolfach Reports Open Problems.'' The record gives proposed year 2016 and does not name a proposer.
\item Supplied source citation: \emph{Surface Bundles} (2017), Oberwolfach Reports, DOI 10.4171/owr/2016/56, pp.\ 3149--3195. \url{https://doi.org/10.4171/owr/2016/56}. The background describes Oberwolfach Reports as documenting mathematical workshops and their open-problem discussions. The supplied literature-source label is ``Oberwolfach Report 56 (2017).''
\item Supplied literature source: van der Geer, ``Counting curves over finite fields.'' \url{https://arxiv.org/abs/1409.6090}.
\item Supplied literature assessment, checked on exactly \texttt{2026-08-21}: the record labels the problem open and states that point counts are computable in low genus and related cohomological data are known in ranges, but that the positive finite limit for the normalized coarse-moduli count is not established; recent asymptotic discussions are described as conjectural or heuristic. This is supplied dated curation derived from the OWR material, not an independently verified literature assessment. The supplied background expressly calls for independent verification of time-sensitive status and cited claims.
\end{itemize}
\end{document}
